% ===========================================================================
\section{Trapping, the cat's eye and the saturation of the shear instability}\label{app:trapping}
% ===========================================================================

Here we derive, step by step, the results that determine the amplitude of the vortices: that the motion of a fluid element in the frame of a vortex is Hamiltonian, \cref{eq:hamiltonian}; that near a critical layer it is the motion of a pendulum, \cref{eq:pendulum}, with the trapping frequency and the half-width of \cref{eq:trapping}; and that an instability of the jets stops growing when these trapped orbits form, which gives the saturation law \cref{eq:saturation}. Only classical mechanics and elementary fluid dynamics are needed.

\subsection{Motion in the frame of the vortex}\label{app:trapping_frame}

Across a strong magnetic field, every particle, whatever its charge and its energy, moves at the \exb{} velocity. In the local coordinates \((x, y)\), radial and binormal, this velocity is derived from the streamfunction \(\Psi\), which is the electrostatic potential in suitable units: \(\dot x = -\partial_y\Psi\) and \(\dot y = \partial_x\Psi\). The flow is everywhere tangent to the contours of \(\Psi\), since \(\dot x\,\partial_x\Psi + \dot y\,\partial_y\Psi = 0\), so each fluid element would stay on one contour if \(\Psi\) did not depend on time. For a vortex that travels on a zonal flow, however,
\begin{equation}
\Psi = \Psi_\zon(x) + \Psi_v(x, y - ct),
\label{eq:G_psi}
\end{equation}
where the zonal part gives the flow \(U(x) = \Psi_\zon'(x)\) in the \(y\) direction, and the part due to the vortex depends on time, because the vortex moves at its phase speed \(c\), but only through the combination \(y - ct\), as long as its shape is steady. Let us, then, measure position from a point that travels with the vortex, \(\xi \equiv y - ct\). Since \(\partial_y = \partial_\xi\) and \(\dot\xi = \dot y - c\), we have \(\dot x = -\partial_\xi\Psi_v\) and \(\dot\xi = U(x) + \partial_x\Psi_v - c\), whose right-hand sides are the derivatives of a single function:
\begin{equation}
\dot x = -\partial_\xi\Ham,
\qquad
\dot\xi = \partial_x\Ham,
\qquad
\Ham(x, \xi) \equiv \Psi_\zon(x) + \Psi_v(x, \xi) - cx.
\label{eq:G_ham}
\end{equation}
These are Hamilton's equations, with \(\xi\) in the role of the coordinate and \(x\) in that of its conjugate momentum; the term \(-cx\) is the streamfunction of the uniform flow \(-c\) that an observer moving with the vortex sees. Since \(\Ham\) does not depend on time explicitly, \(\rmd\Ham/\rmd t = \partial_x\Ham\,\dot x + \partial_\xi\Ham\,\dot\xi = 0\): each fluid element stays on one contour of \(\Ham\), viz., on a streamline of the flow as it is seen from the vortex. Equation \cref{eq:G_ham} is \cref{eq:hamiltonian} without its drift term, \(\vda x\), which arises in the same way from the magnetic drift and is negligible for the motion (\cref{sec:geometry}).

\subsection{The pendulum and the cat's eye}\label{app:trapping_pendulum}

Without the vortex, \cref{eq:G_ham} gives \(\dot\xi = U(x) - c\): the fluid that flows faster than the vortex overtakes it, the fluid that flows more slowly falls behind, and the fluid on the line \(x = x_r\), where \(U(x_r) = c\), stays level with it. This line is the critical layer. Let us expand \(\Ham\) about it. With \(X \equiv x - x_r\) and the shear \(S \equiv U'(x_r)\), Taylor's theorem gives \(\Psi_\zon(x) - cx = \text{const} + [U(x_r) - c]X + SX^2/2 + \dots\), in which the term linear in \(X\) vanishes by the definition of \(x_r\). If the vortex varies little in \(x\) over the region of interest, we may also set \(\Psi_v(x, \xi) \approx \Psi_v(x_r, \xi) = -A\cos k\xi\), where \(k\) is the wavenumber of the vortex, \(A > 0\) is its amplitude at the critical layer, and the origin of \(\xi\) has been put at a minimum of \(\Psi_v\). Dropping the constant, we obtain \cref{eq:pendulum} and its equations of motion,
\begin{equation}
\Ham = \frac{S}{2}X^2 - A\cos k\xi,
\qquad
\dot\xi = SX,
\qquad
\dot X = -kA\sin k\xi.
\label{eq:G_pendulum}
\end{equation}
The first equation of motion says that the shear carries an element at the distance \(X\) from the critical layer past the vortex at the relative speed \(SX\); the second, that the vortex pushes the element across the layer. Differentiating the first and substituting the second,
\begin{equation}
\frac{\rmd^2(k\xi)}{\rmd t^2} = -k^2SA\,\sin k\xi,
\label{eq:G_pend_eq}
\end{equation}
which is the equation of a pendulum whose angle is \(k\xi\).

Let us take \(S > 0\) first. The point \(k\xi = 0\), \(X = 0\) is the stable equilibrium of the pendulum, and is called the O-point; for small oscillations about it, \(\sin k\xi\approx k\xi\) in \cref{eq:G_pend_eq}, and the motion is harmonic with the frequency \(\omtr\) given below. The points \(k\xi = \pm\pi\), \(X = 0\) are the unstable equilibrium, the pendulum balanced upside down, and are called X-points. The other orbits follow from the conservation of \(\Ham\), whose value is \(-A\) at the O-point and \(+A\) at the X-points. An orbit with \(\Ham < A\) cannot reach \(k\xi = \pm\pi\), where \cref{eq:G_pendulum} would require \(SX^2/2 = \Ham - A < 0\); it goes round the O-point on a closed curve and is said to be trapped. An orbit with \(\Ham > A\) has \(SX^2/2 = \Ham + A\cos k\xi \geq \Ham - A > 0\), so \(X\) never changes sign: the element stays on one side of the critical layer and overtakes the vortex, or is overtaken by it, for ever. The boundary between the two families is the separatrix, \(\Ham = A\), on which \(SX^2/2 = A(1 + \cos k\xi) = 2A\cos^2(k\xi/2)\). Its two arcs, \(X = \pm w\cos(k\xi/2)\), meet at the X-points and enclose a lens-shaped region of trapped fluid, the cat's eye [see \cref{fig:cartoon}(a) and \cref{fig:geometry}(b)], where
\begin{equation}
\omtr = |k|\sqrt{|S|A},
\qquad
w = 2\sqrt{\frac{A}{|S|}}
\label{eq:G_trapping}
\end{equation}
are the trapping frequency and the half-width of the eye at its widest, i.e., \cref{eq:trapping}. Both scale as \(A^{1/2}\). The approximations made in \cref{eq:G_pendulum} hold if \(\Psi_v\) and the shear change little over the distance \(w\).

Two further properties of these orbits will be needed. Firstly, the trapped orbits do not all have the same period: it is \(2\pi/\omtr\) at the centre of the eye and increases outwards, becoming infinite on the separatrix, because a pendulum released from its inverted position takes for ever to leave it. Any pattern carried by the fluid inside the eye is therefore wound into an ever finer spiral. Secondly, if \(S < 0\), the right-hand side of \cref{eq:G_pend_eq} changes sign, so the O-point is at \(k\xi = \pi\) and the X-points are at \(k\xi = 0\), while \cref{eq:G_trapping} holds as written. A jet has \(S > 0\) on one of its flanks and \(S < 0\) on the other, so a vortex whose phase is the same across the jet has its two rows of eyes displaced from each other by half a wavelength, as in \cref{fig:cartoon}(a).

When the variation of the vortex in \(x\) cannot be neglected, the same reasoning applies to the full \(\Ham\). An equilibrium is a point where \(\partial_x\Ham = \partial_\xi\Ham = 0\). Expanding \cref{eq:G_ham} to first order in the displacements \((\delta x, \delta\xi)\) from it gives \(\delta\dot x = -\Ham_{x\xi}\,\delta x - \Ham_{\xi\xi}\,\delta\xi\) and \(\delta\dot\xi = \Ham_{xx}\,\delta x + \Ham_{x\xi}\,\delta\xi\), where the subscripts denote second derivatives at the equilibrium, and solutions proportional to \(\rme^{\lambda t}\) exist if \(\lambda^2 = \Ham_{x\xi}^2 - \Ham_{xx}\Ham_{\xi\xi}\). If this is negative, the neighbouring orbits circle the equilibrium, which is an O-point, at the frequency \(\omtr\) given by \(\omtr^2 = \Ham_{xx}\Ham_{\xi\xi} - \Ham_{x\xi}^2\); if it is positive, they leave it exponentially, and it is an X-point. For \cref{eq:G_pendulum} at its O-point, \(\Ham_{xx} = S\), \(\Ham_{\xi\xi} = k^2A\) and \(\Ham_{x\xi} = 0\), and \cref{eq:G_trapping} is recovered.

\subsection{Saturation of the instability}\label{app:trapping_saturation}

Suppose now that the jets are unstable, i.e., that Rayleigh's equation \cref{eq:rayleigh} has a mode that grows at the rate \(\gamma_\text{R}\) (\cref{sec:euler}). The linear calculation that gives \(\gamma_\text{R}\) follows each fluid element along the path that it would take without the wave, viz., \(X = \text{const}\) and \(\xi = \xi_0 + SXt\) by \cref{eq:G_pendulum}, and computes the small radial displacement \(\delta X\) that the wave adds to it. For a wave of amplitude \(A(t)\propto\rme^{\gamma_\text{R}t}\), integrating the last of \cref{eq:G_pendulum} along this path from the distant past gives
\begin{equation}
\delta X = -k\int_{-\infty}^{t}A(t')\sin\!\left(k\xi_0 + kSXt'\right)\rmd t'
= -kA(t)\operatorname{Im}\frac{\rme^{\rmi k\xi}}{\gamma_\text{R} + \rmi kSX},
\label{eq:G_displacement}
\end{equation}
whose magnitude is at most \(kA/(\gamma_\text{R}^2 + k^2S^2X^2)^{1/2}\). The meaning of \cref{eq:G_displacement} is as follows. An element at the distance \(X\) from the critical layer sees the wave at the frequency \(kSX\), Doppler shifted by its motion past the vortex. If \(|kSX|\gg\gamma_\text{R}\), the push that it receives reverses many times while the wave grows by a factor of \(\rme\), and its effect averages out. If \(|kSX|\lesssim\gamma_\text{R}\), the push has the same sign for a whole growth time. It is these elements, displaced steadily across the gradient of the zonal vorticity, that pass energy from the jet to the wave [cf.\ \cref{eq:gamma_rayleigh}]. They occupy a resonant layer about the critical layer, of half-width \(X_\gamma \equiv \gamma_\text{R}/|k||S|\).

The true paths, however, are those of \cref{eq:G_pendulum}, and the unperturbed path is a good approximation to them only as long as the wave cannot turn an element round within one growth time, i.e., as long as \(\omtr\ll\gamma_\text{R}\). By \cref{eq:G_trapping}, the same condition can be written as
\begin{equation}
\frac{w}{X_\gamma} = \frac{2\omtr}{\gamma_\text{R}} \ll 1,
\label{eq:G_ratio}
\end{equation}
viz., the cat's eye must be a thin sliver inside the resonant layer. As the wave grows, \(\omtr\propto A^{1/2}\) grows with it, until \(\omtr\) is comparable with \(\gamma_\text{R}\). The eye has then swallowed the resonant layer, and the elements that were feeding the wave go round the eye in about the time that the wave needs to grow. An element that goes round gives the wave nothing on average: on one half of its orbit it moves across the layer one way, and on the other half it comes back. Moreover, since the orbits have different periods, the vorticity inside the eye is wound up and, after a few turns, is uniform on average along each closed streamline, so the gradient that drove the instability has been flattened where it mattered. The growth stops. This is the argument of \citet{ONeil1965} for a plasma wave and the particles that resonate with it; for a shear flow, it is the theory of the nonlinear critical layer \citep{Stewartson1978, WarnWarn1978}.

The wave therefore saturates at the amplitude for which
\begin{equation}
\omtr = \alpha\,\gamma_\text{R},
\qquad\text{i.e.,}\qquad
A = \frac{\alpha^2\gamma_\text{R}^2}{k^2|S|},
\label{eq:G_saturation}
\end{equation}
where \(\alpha\) is a pure number and \cref{eq:G_trapping} has been used. This is \cref{eq:saturation}. The argument fixes how the amplitude scales with the growth rate, the wavenumber and the shear, which is what can be tested, but not the value of \(\alpha\): `comparable' covers anything from a fraction of a trapping period, \(2\pi/\omtr\), to several periods, and the answer depends on how the vorticity is distributed near the layer. We do not derive \(\alpha\); it is measured in \cref{fig:saturation}.

Finally, let us explain why the amplitude continues to obey \cref{eq:G_saturation} while the jets change. Here \(\gamma_\text{R}\) is the growth rate on the jets as they are at a given time, so \cref{eq:G_saturation} relates two quantities that can be measured independently. The trapping time \(1/\omtr\) and the growth time \(1/\gamma_\text{R}\) are comparable by \cref{eq:G_saturation}, and both are much shorter than the time over which the jets, and hence \(\gamma_\text{R}\), change. If the amplitude is below the level \cref{eq:G_saturation} of the current jets, the eye is narrower than the resonant layer, and the wave grows again, at a rate of order \(\gamma_\text{R}\), until the elements that feed it are trapped once more; this takes a few growth times, an instant on the time scale of the jets, and it is what happens when the vortices are reduced artificially [see \cref{fig:euler}(b)]. If, on the other hand, \(\gamma_\text{R}\) falls, the level that trapping can sustain falls with it, and the weak damping of the vortices (\cref{sec:dissipation}), which acts all the time, brings the amplitude down to that level, provided that the jets change more slowly than the vortices are damped. In other words, the amplitude follows \(\gamma_\text{R}(t)\) adiabatically, as in \cref{fig:saturation}(a), until \(\gamma_\text{R}\) vanishes, the jets are stable, and nothing sustains the vortices.
